\section{Pole coefficients and zero zeta functions}
\subsection{Nearest-pole estimates retain a local analytic contract}
For a positive integral reciprocal order, a simple denominator zero produces a pole of finite order. The source constructs principal parts, proves their uniqueness, removes them, and bounds the Taylor coefficients of the analytic remainder. This produces the nearest-circle coefficient expansion with a controlled error.

The generic coefficient bound below states its local principal-part assumption explicitly. Later source results construct the required principal parts for the manuscript specialization.
\excerpt{NearestPoleExpansion}{nearest_pole}
The finite set $S$ contains the zeros inside the larger disk. The inner radius $R$ is strictly smaller than the outer radius $R_2$. The assumption \code{hlocal} states that after subtracting the proposed principal part, the function is locally bounded near each selected pole. These are the hypotheses needed by the removable-singularity and Taylor-bound route.

The conclusion is in the manuscript's block normalization. Its leading sum and error both carry $(Nn)!$. The factor $\rho^{-Nn}$ arises from extracting the original-variable coefficient at degree $Nn$. The binomial coefficient records the higher pole order. Omitting any of these factors changes the asymptotic statement.

\subsection{Principal parts are constructed before the final specialization}
Rows 133--138 provide coefficient bounds, principal-part coefficients, pole removal, Taylor recovery, simple-pole expansions, and uniqueness. Row 139 supplies the full nearest-pole expansion. Row 140 isolates the leading term under its stated dominance hypotheses; row 141 reduces symmetry orbits on the nearest circle; row 142 gives the secant specialization.

This route is conditional in the mathematically ordinary sense that the theorem retains its hypotheses on reciprocal order, selected zeros, radii, and simplicity where used. Formal completion proves those scoped statements. It does not delete the assumptions or claim that every parameter and every zero configuration satisfies them.

\subsection{Absolute zero zeta and branch-dependent zeta differ}
The absolute zero zeta uses terms of the form
\[
 m(\rho)\exp\bigl(-s\log |\rho|\bigr),
\]
where $m(\rho)$ is analytic multiplicity and the real logarithm of the norm is unambiguous. The reduced denominator has sharp absolute-convergence half-plane $\operatorname{Re}s>1/N$; the original denominator has $\operatorname{Re}s>1$. Local uniform convergence and holomorphy are proved, as is failure of absolute convergence on and to the left of the respective boundaries.

A branch-dependent symmetry-orbit zeta instead assigns complex logarithms. Its terms carry argument information. The source requires compatible logarithm choices and records the phase factors in its continuation. These are related constructions, not interchangeable notations for one function.

\subsection{Continuation is an explicit model plus an entire correction}
Exponential displacement of the positive zero sequence from a lattice makes the difference of the two zeta sums entire. The source proves a locally uniform majorant on bounded sets of the parameter. It then constructs the lattice correction, finite missing lattice prefix, and finite exceptional-zero contribution.
\excerpt{ZeroZetaContinuation}{zeta_correction}
The finite exception term retains multiplicity. The correction and prefix use the scaled parameter $Ns$ because reduced roots are $N$th powers of the positive original roots. The equality with the Hurwitz lattice model is proved in the initial convergence domain:
\excerpt{ZeroZetaContinuation}{zeta_identity}
The continuation is then defined by the meromorphic model plus its entire correction.
\excerpt{ZeroZetaContinuation}{zeta_meromorphic}
This theorem is short because its two analytic ingredients have already been proved. The continuation is tied to the original sum by the preceding equality. The source also proves holomorphy off the pole, exact residues, and meromorphic order $-1$ at the nonzero-residue pole. It does not merely define a meromorphic function and call it a continuation.

\subsection{Derivative weights and boundary orders}
Derivative-weighted sums use complex derivative values, not only their moduli. A lower bound on the derivative at sufficiently large zeros gives exponential decay of its inverse powers. For $N\geq3$ and positive integral derivative weight, this dominates every fixed polynomial weight and yields an entire sum, assuming logarithm arguments remain in a fixed bounded interval.

The order-two case is separate: the derivative magnitude does not provide the same exponential gain, and the theorem retains a right half-plane of convergence. Rows 85--87 state and prove this distinction. Exact order-two and order-four zeta models are separately identified in rows 94--95.
